% TOP_PROBLEM: {"id": 30007336, "title": "Smooth anisotropic conductivity uniqueness from a single proper patch", "category_id": 9, "category": {"id": 9, "name": "pde", "display_name": "Partial Differential Equations", "description": "PDEs and their applications in physics and geometry.", "slug": "pde", "order_index": 9, "created_at": "2026-07-31T15:26:25.671Z"}, "status": "open"}
% FIRST_POSTED: 2026-10-06
% ATTEMPT_STATUS: solved
% OPENAI_CLAIM: {"schema_version":1,"source_repo":"https://github.com/openai/math","source_commit":"adc7f1241b42e322a6451854ab7e4b4c146bf78a","release_date":"2026-10-06","match":"full","resolution":"proved","independently_reviewed":false,"second_pass":"agreed","families":[{"family":"365","title":"Joint metric and connection recovery from one boundary patch","subject":"Partial differential equations","manuscripts":[{"title":"Smooth anisotropic uniqueness in the Calderón problem from one boundary patch","pdf_path":"preprints/Smooth-Anisotropic-Uniqueness-in-the-Calderon-Problem-from-One-Boundary-Patch-September-24-2026/paper.pdf","readme_path":"preprints/Smooth-Anisotropic-Uniqueness-in-the-Calderon-Problem-from-One-Boundary-Patch-September-24-2026/README.md","date":"2026-09-24","theorem_ref":"main:full (sections/01-introduction.tex:84); main:patch (sections/01-introduction.tex:54)","statement_sources":[{"path":"preprints/Smooth-Anisotropic-Uniqueness-in-the-Calderon-Problem-from-One-Boundary-Patch-September-24-2026/build/sections/01-introduction.tex","line":84,"label":"main:full"},{"path":"preprints/Smooth-Anisotropic-Uniqueness-in-the-Calderon-Problem-from-One-Boundary-Patch-September-24-2026/build/sections/01-introduction.tex","line":54,"label":"main:patch"}]}],"lean":{"doc_path":"lean/docs/365.md","comparators":[],"declarations":[],"covers_main_theorem":false},"reasoning_trace":null}]}
\documentclass[11pt]{article}
\usepackage[margin=25mm]{geometry}
\usepackage{fontspec}
\setmainfont{Latin Modern Roman}
\usepackage{amsmath,amssymb,mathtools}
\usepackage{unicode-math}
\setmathfont{Latin Modern Math}
\usepackage{enumitem}
\usepackage{xurl,hyperref,bookmark}
\hypersetup{colorlinks=true,urlcolor=blue}
\setcounter{secnumdepth}{0}
\setlength{\parindent}{0pt}
\setlength{\parskip}{5pt}
\setlength{\emergencystretch}{3em}
\begin{document}
\section*{Smooth anisotropic conductivity uniqueness from a single proper patch}

\subsection{OpenAI's claim}
\textit{One-patch anisotropic uniqueness}: The paper claims determination of a smooth Riemannian metric in dimension at least three from its zero-frequency energy form on any one open boundary patch. The only ambiguity is a diffeomorphism equal to the identity on that patch.

\subsection{Mathematical statement of the claimed result}
\paragraph{Mathematical claim 1.} Fix a dimension $n\ge3$. Assume the smooth $n$-manifold $M$ is compact and connected, with smooth nonempty boundary. For a proper relatively open observation patch $\Gamma\ne\varnothing$, the claim asserts \[\Lambda_{g_1,\Gamma}=\Lambda_{g_2,\Gamma}\quad\Longrightarrow\quad g_2=\Phi^*g_1,\qquad\Phi|_\Gamma=\operatorname{id},\] for some diffeomorphism smooth through the boundary. Both metrics are smooth, and the local maps are the harmonic Dirichlet energy forms for data supported in $\Gamma$. Equality for all full-boundary data gives a diffeomorphism fixing the entire boundary. This concerns smooth anisotropic metrics, not arbitrary bounded measurable scalar conductivities.

\subsection{Outline of the argument}
\textit{One-patch anisotropic uniqueness}: Boundary jets and a common exterior cap identify actual Green kernels in an initial region. Remote source solutions produce harmonic coordinates and enough second-order jets. A mixed Green-kernel continuation theorem extends metric equality across a smooth frontier, so the maximal paired coordinate region cannot stop in the interior.

\subsection{Scope relative to this problem}
\textit{One-patch anisotropic uniqueness}: The metrics are smooth and the dimension is at least three; the measurements can occupy a proper open part of the boundary. Determination is up to a smooth diffeomorphism fixing that measured patch, preserving the natural anisotropic ambiguity.

With the hypotheses and conventions above, the claimed result addresses the whole mathematical assertion on this page.

\subsection{Formal verification}
Family 365 describes Lean coverage for the measurable scalar-conductivity obstruction; smooth anisotropic and coupled connection recovery are separate.

\textit{One-patch anisotropic uniqueness}: No manuscript-specific Comparator is listed in the pinned inventory.

The documented selection does not cover the entire manuscript assertion relevant here.

This site has not run the Lean development. These links describe OpenAI’s selected formal statements and should be read with their stated scope.

\subsection{Status caveat}
OpenAI reports that an internal model produced the manuscripts. Their arguments have not been independently verified as proofs by this site. Any portion outside the documented formal scope may contain mistakes, and the catalogue label reports the announced mathematical claim.

\subsection{Sources}
\begin{itemize}
\item OpenAI, \emph{Smooth anisotropic uniqueness in the Calderón problem from one boundary patch}, 2026-09-24; citation key \texttt{OAI:Smooth-Anisotropic-Uniqueness-in-the-Calderon-Problem-from-One-Boundary-Patch-September-24-2026}. \href{https://github.com/openai/math/blob/adc7f1241b42e322a6451854ab7e4b4c146bf78a/preprints/Smooth-Anisotropic-Uniqueness-in-the-Calderon-Problem-from-One-Boundary-Patch-September-24-2026/paper.pdf}{PDF (GitHub)}; \href{https://github.com/openai/math/raw/adc7f1241b42e322a6451854ab7e4b4c146bf78a/preprints/Smooth-Anisotropic-Uniqueness-in-the-Calderon-Problem-from-One-Boundary-Patch-September-24-2026/paper.pdf}{PDF (direct)}; \href{https://github.com/openai/math/blob/adc7f1241b42e322a6451854ab7e4b4c146bf78a/preprints/Smooth-Anisotropic-Uniqueness-in-the-Calderon-Problem-from-One-Boundary-Patch-September-24-2026/README.md}{Manuscript page and citation}; \href{https://github.com/openai/math/blob/main/preprints/Smooth-Anisotropic-Uniqueness-in-the-Calderon-Problem-from-One-Boundary-Patch-September-24-2026/paper.pdf}{Latest version}; Statement source: \href{https://github.com/openai/math/blob/adc7f1241b42e322a6451854ab7e4b4c146bf78a/preprints/Smooth-Anisotropic-Uniqueness-in-the-Calderon-Problem-from-One-Boundary-Patch-September-24-2026/build/sections/01-introduction.tex\#L84}{main:full}; Statement source: \href{https://github.com/openai/math/blob/adc7f1241b42e322a6451854ab7e4b4c146bf78a/preprints/Smooth-Anisotropic-Uniqueness-in-the-Calderon-Problem-from-One-Boundary-Patch-September-24-2026/build/sections/01-introduction.tex\#L54}{main:patch}.
\item Family 365: \href{https://github.com/openai/math/blob/adc7f1241b42e322a6451854ab7e4b4c146bf78a/lean/docs/365.md}{Lean scope}.
\item \href{https://github.com/openai/math/blob/adc7f1241b42e322a6451854ab7e4b4c146bf78a/CONTENTS.md}{OpenAI catalogue entry}.
\end{itemize}
\end{document}
